\documentclass[11pt]{article}

\usepackage[a4paper,margin=28mm]{geometry}
\usepackage[T1]{fontenc}
\usepackage[utf8]{inputenc}
\usepackage{lmodern}
\usepackage{microtype}
\usepackage{amsmath,amssymb,amsthm,mathtools}
\usepackage{booktabs,longtable,array}
\usepackage{enumitem}
\usepackage{xcolor}
\usepackage[hidelinks]{hyperref}
\usepackage{url}

\hypersetup{
  pdftitle={Fifteen lonely runners},
  pdfauthor={Jaan Allikvere}
}

\newtheorem{theorem}{Theorem}[section]
\newtheorem{proposition}[theorem]{Proposition}
\newtheorem{lemma}[theorem]{Lemma}
\newtheorem{corollary}[theorem]{Corollary}
\theoremstyle{definition}
\newtheorem{definition}[theorem]{Definition}
\theoremstyle{remark}
\newtheorem{remark}[theorem]{Remark}

\newcommand{\Z}{\mathbb{Z}}
\newcommand{\R}{\mathbb{R}}
\newcommand{\vct}[1]{\mathbf{#1}}
\newcommand{\nearest}[1]{\left\lVert #1\right\rVert}
\newcommand{\Enorm}[1]{\left\lVert #1\right\rVert_{E}}
\DeclareMathOperator{\lcm}{lcm}
\DeclareMathOperator{\spn}{span}
\DeclareMathOperator{\covol}{covol}

% ---- verification constants (filled from the certified data) ----
% \GateCount, \GateMassLB (rounded DOWN), \NoNearGcdCount, \NearGcdCount,
% \CostExponent, \TotalCpuHours, \MassMargin are generated from the certified
% data by inputs/make_tables.py.
\newcommand{\GateCount}{71}
\newcommand{\GateMassLB}{408.8233}
\newcommand{\NoNearGcdCount}{52}
\newcommand{\NearGcdCount}{19}
\newcommand{\CostExponent}{5.9}
\newcommand{\CpuExponent}{6.2}
\newcommand{\TotalCpuHours}{2,137}
\newcommand{\TotalCpuYears}{0.24}
\newcommand{\MassMargin}{6.84}

% Thresholds are rounded UP, so every strict inequality below survives the rounding.
\newcommand{\RawThreshUB}{414.7794}       % 14 log(A_14/28) rounded up
\newcommand{\DivisorLogLB}{12.7948}       % log 360360 rounded down
\newcommand{\TargetUB}{401.9846}          % RawThreshUB - DivisorLogLB, rounded up
\newcommand{\OldThresh}{810.0740}         % log B_14 of the MSS/ST bound, rounded up

\title{Fifteen lonely runners and a lattice flag bound for finite checking}
\author{Jaan Allikvere\\
\small Independent researcher, Tallinn, Estonia\\
\small ORCID: \href{https://orcid.org/0009-0003-5228-7015}{0009-0003-5228-7015}}
\date{September 2026 --- draft}

\begin{document}
\maketitle

\begin{abstract}
We prove the Lonely Runner Conjecture for fifteen runners.  The proof
combines a new finite-checking bound with exhaustive computations for
\GateCount{} primes.  Motivated by Giri and Kravitz, we study the
projection of $\Z^n$ onto the orthogonal complement of a primitive speed
tuple.  We extend their two-dimensional density argument along a full
flag of rational subtori, and nearest-plane
rounding turns the resulting estimates into simultaneous inequalities for
a Korkine--Zolotarev basis.  An exact covolume formula then bounds the
product of the speeds.  For $n=14$, the logarithm of this bound is
$414.779\ldots$, compared with $810.074\ldots$ for the bound used in the
preceding computer-assisted proofs.  A simple forced-divisor argument
reduces the logarithmic sum required from the verified primes to
$\TargetUB$.

The modular computation begins with a complete two-branch covering search.
It treats the terminal level $15=3\cdot5$ through its factors and uses a
few-exception shift lemma for the residual classes left by the modular
framework.  The verified primes have
logarithmic sum greater than $\GateMassLB$, which exceeds the required
threshold.  Code, certificates, and exact checks of the constants accompany
the proof.
\end{abstract}

\medskip
\noindent\textbf{Keywords.}
Lonely Runner Conjecture; finite checking; Korkine--Zolotarev reduction;
nearest-plane rounding; computer-assisted proof; exhaustive computation.

\section{Introduction}

For a real number $x$, let $\nearest{x}$ denote its distance to the nearest
integer.  The Lonely Runner Conjecture $LRC(k)$ asserts that for every
$\vct{v}=(v_1,\ldots,v_k)\in(\Z\setminus\{0\})^k$ there is a real time $t$
with
\begin{equation}\label{eq:lrc}
  \nearest{t v_i}\geq \frac{1}{k+1}\qquad(1\leq i\leq k).
\end{equation}
This is the statement for $k+1$ runners after one runner is made
stationary; since $\nearest{-x}=\nearest{x}$, the speeds may be taken
positive.  We call $\vct{v}$ a \emph{speed tuple}.  If such a time $t$
exists, the tuple has the \emph{LR property}, and $t$ is a \emph{witness}.

The problem goes back to work of Wills in the 1960s~\cite{Wills1967}.
Rosenfeld proved $LRC(7)$~\cite{Rosenfeld25}; Trakulthongchai proved
$LRC(8)$ and $LRC(9)$~\cite{Trakulthongchai26}; Sungkawichai and
Trakulthongchai proved $LRC(k)$ for $10\leq k\leq12$~\cite{ST26}; and the
author proved $LRC(13)$~\cite{Allikvere26}.

\begin{theorem}\label{thm:main}
The Lonely Runner Conjecture holds for fifteen runners.
\end{theorem}

The recent cases use the same prime-divisibility strategy.  A
finite-checking theorem bounds the product of the speeds in a primitive
counterexample.  For each verified prime $p$, a modular computation shows
that $p$ divides this product.  Once the product of the verified primes
exceeds the finite-checking bound, no counterexample can exist.

For $LRC(14)$, however, the earlier bound would require primes extending to
about $970$.  Our measurements indicate that the cost at a prime $p$ grows
roughly as $p^6$, so this direct continuation would be expensive.  The
proof therefore has two parts.  First, we improve the finite-checking bound
by studying the projected lattice in a norm adapted to the speed tuple.
Lower-dimensional cases of the conjecture give an inequality for every
prefix of a Korkine--Zolotarev basis, rather than only the top-dimensional
inequality used previously.  For fourteen speeds this gives
\[
  \log(v_1\cdots v_{14})<414.779\ldots,
\]
compared with $810.074\ldots$ from the earlier bound.  The elementary
divisibility $360360=\lcm(2,\ldots,15)\mid v_1\cdots v_{14}$ reduces the
logarithmic sum required from the verified primes to $\TargetUB$.

Second, we verify the required divisibility for \GateCount{} primes.  A
complete two-branch covering search generates the level-$1$ improper
tuples.  Successive binary lifting removes most of them.  At the terminal
level we factor $15=3\cdot5$, and a few-exception shift lemma supplies the
remaining divisibility argument when the covering number is low.  The logarithms of
the verified primes sum to more than $\GateMassLB$.

\medskip
\noindent\textbf{Outline of the proof.}
Section~\ref{sec:framework} recalls the modular framework and proves the
prime-divisibility criterion used here.  Section~\ref{sec:flag} proves the
lattice flag bound and the forced-divisor lemma.  Section~\ref{sec:pipeline}
gives the verification for a fixed prime, including completeness proofs for
the new steps.  Section~\ref{sec:gateset} lists the verified primes and
deduces Theorem~\ref{thm:main}.  The final main-text sections discuss
reproducibility, the surviving orbits, and the next case.  The appendices
give two complete examples, the detailed verification record,
computational costs, and the source check for the $LRC(9)$ input.

\medskip
\noindent\textbf{What is new.}
The analytic contribution is a finite-checking theorem based on the full
flag of prefix inequalities and an exact covolume in a speed-adapted
ellipsoidal norm.  On the computational side, the two-branch generator of
\cite{Allikvere26} is extended by a proved general variant, while the
factorisation of level $15$ and the few-exception shift lemma replace a
prohibitively large direct search.  Each verified prime has a certificate,
and separate scripts check the certificates and recompute the analytic
constants.  The final prime-product deduction uses the framework
of~\cite{ST26} and takes the known cases through $LRC(13)$ as input.

\subsection*{Relation to earlier work}

Two earlier finite-checking arguments are especially relevant.
Malikiosis, Santos, and Schymura~\cite{MSS25} obtain the product bound used
in the preceding computer-assisted proofs; Blanco, Criado, and
Santos~\cite{BCS26} recover the same bound in a broader setting.  Giri and
Kravitz~\cite{GK26} instead use density in a two-dimensional rational
subtorus, where the $LRC(n-1)$ case---equivalently, the case of $n$
runners---applies.  Our argument follows this geometric approach along a
complete flag of rational subtori.  Nearest-plane rounding turns the
density statements into prefix inequalities for a Korkine--Zolotarev
basis, while the adapted ellipsoid gives the covolume exactly.  The full
comparison appears after Theorem~\ref{thm:flag}.

\section{Proof strategy and the prime-divisibility criterion}\label{sec:framework}

We recall the definitions of Sungkawichai and
Trakulthongchai~\cite{ST26} in the form needed here, then isolate the
prime-divisibility step used in the final argument.

Let $p$ be prime and $l$ a positive integer, and write
$\Z_{p,l}=\Z_{pl}\setminus p\Z_l$, so that a vector in $\Z_{p,l}^k$ has no
coordinate divisible by $p$.  A vector $\vct{v}\in\Z_{p,l}^k$ is
\emph{$(k,p,l)$-proper} if either (a) for some $i$,
$\gcd(l,v_1,\ldots,\widehat{v_i},\ldots,v_k)>1$, or (b) for some
$t\in(lp)^{-1}\Z$, $\nearest{tv_i}\geq1/(k+1)$ for every $i$.  The set of
improper vectors is $I(k,p,l)$.  A vector $\vct{v}\in\Z_{p,1}^k$ is
\emph{eventually $(k,p)$-proper} if some level $l$ has no improper lift of
$\vct{v}$, where a lift to level $l$ is a vector
$\vct{w}\equiv\vct{v}\pmod p$ with coordinates in $\Z_{pl}$; the set of
vectors that are not eventually proper is $J(k,p)$.  By Lemma~2.4
of~\cite{ST26}, $J(k,p)=\varnothing$ if and only if $I(k,p,l)=\varnothing$
for some $l$.  The equivalence uses the following monotonicity: a witness,
or a divisor from the gcd condition, remains valid at every multiple of
the level.  By Proposition~5.1 of~\cite{ST26}, permutations,
sign changes of coordinates, and multiplication by a unit modulo $p$
preserve eventual properness (for the unit case see also Remark~2.3
of~\cite{Allikvere26}).

For an integer $c\geq2$, the $c$-lifts of $\vct{v}\in\Z_{p,l}^k$ to level
$cl$ are the $c^k$ vectors $w_i=v_i+a_ilp$, $a_i\in\{0,\ldots,c-1\}$; these
vectors form the full lift fiber of $\vct{v}$.  For a level-$1$ vector $r$
and a level $l$ write
$F_l(r)=\{\vct{w}\in I(k,p,l):\vct{w}\equiv r\bmod p\}$.  The computation
carries exactly this improper part of the fiber for every generated $r$
and every level it reaches.  This is true at level $1$ and propagates
through each lifting stage by monotonicity, since every element of
$F_{cl}(r)$ reduces to an element of $F_l(r)$.  If $F_l(r)=\varnothing$ at
some level, $r$ is eventually proper, and so is its unit orbit.

\begin{lemma}[Prime-divisibility criterion]\label{lem:gate}
Let $k\geq3$, assume $LRC(m)$ for all $m<k$, and let $\vct{v}\in\Z_{>0}^k$
be a primitive counterexample to $LRC(k)$ with pairwise distinct
coordinates.  Let $p$ be a prime.
\begin{enumerate}[leftmargin=*,label=(\roman*)]
\item If $J(k,p)=\varnothing$, then $p\mid v_1\cdots v_k$.
\item More generally, suppose that for every $\vct{u}\in\Z_{p,1}^k$
  either $\vct{u}$ is eventually proper, or there is a level $l$ such
  that every lift $\vct{w}$ of $\vct{u}$ to level $l$ is proper or
  satisfies the following: every integer vector congruent to $\vct{w}$
  modulo $lp$ with pairwise distinct positive coordinates has the
  lonely-runner property~\eqref{eq:lrc}.  Then $p\mid v_1\cdots v_k$.
\end{enumerate}
\end{lemma}

\begin{proof}
Part (i) is Lemma~2.2 of~\cite{ST26} combined with Lemma~2.4 there.  We
recall the mechanism, since (ii) extends it.  Suppose $p\nmid v_i$ for all
$i$ and let $\vct{u}$ be the reduction of $\vct{v}$ modulo $p$.  Take a
level $l$ as in the hypothesis and let $\vct{w}$ be the reduction of
$\vct{v}$ modulo $lp$, a lift of $\vct{u}$.  If $\vct{w}$ has a witness
$t\in(lp)^{-1}\Z$, then $\nearest{tv_i}=\nearest{tw_i}\geq1/(k+1)$ for
all $i$, so $\vct{v}$ is not a counterexample.  If
$g=\gcd(l,w_j:j\neq i)>1$, then $g$ divides $v_j$ for all $j\neq i$, and
$g\nmid v_i$ because $\vct{v}$ is primitive.  The $k-1$ speeds $v_j$,
$j\neq i$, take $m\leq k-1$ distinct values, so $LRC(m)$ gives $t_0$ with
$\nearest{t_0v_j}\geq1/(m+1)\geq1/k$ for all $j\neq i$.  The shifted
times $t_0+q/g$, $q=0,\ldots,g-1$, leave every $t_0v_j$ ($j\neq i$) fixed
modulo one and move $t_0v_i$ through the $g'$ equally spaced points of a
grid with $g'=g/\gcd(g,v_i)\geq2$.  Such a grid cannot lie inside an open
arc of length $2/(k+1)\leq\frac12$: the shortest closed arc containing the
grid has length $1-1/g'\geq\frac12$, and equality is harmless because the
bad arc is open.  Thus some shift has
$\nearest{(t_0+q/g)v_i}\geq1/(k+1)$; at that shift every coordinate is at
distance $\geq1/(k+1)$ and $\vct{v}$ is not a counterexample.  In case
(ii), the third alternative applies directly to
$\vct{v}$ itself, which is congruent to $\vct{w}$ modulo $lp$ and has
pairwise distinct positive coordinates.  In every case $\vct{v}$ has the
lonely-runner property, contradicting the assumption; so some $v_i$ is
divisible by $p$.
\end{proof}

Part~(i) handles primes with $J(14,p)=\varnothing$; for the others,
part~(ii) and Lemma~\ref{lem:neargcd} handle the residue classes left after
lifting.  Both give $p\mid v_1\cdots v_k$.

\begin{remark}[Unit orbits]
If $\vct{v}\equiv a r\pmod p$ with $a\in\Z_p^\times$, the Chinese
remainder theorem gives $b>0$ with $b\equiv a^{-1}\pmod p$ and
$b\equiv1\pmod l$, since $\gcd(l,p)=1$.  Then $(b\vct{v}\bmod lp)$ lifts
$r$, while $b\vct{v}$ remains positive and pairwise distinct, has the same
divisibility by divisors of $l$, and is a counterexample if and only if
$\vct{v}$ is.  Thus one orbit representative suffices.
\end{remark}

\section{The lattice flag bound and the forced divisor}\label{sec:flag}

Throughout this section $n\geq2$, $\vct{v}=(v_1,\ldots,v_n)\in\Z_{>0}^n$ is
primitive ($\gcd(v_1,\ldots,v_n)=1$), and ``counterexample'' means: for
every real $t$ some $i$ has $\nearest{tv_i}<1/(n+1)$.  We write
\[
  S=\sum_iv_i,\quad x_i=\frac{v_i}{S},\quad
  u_i=\frac{v_i}{\max_jv_j},\quad
  q_i=1+2u_i-u_i^2\in(1,2],\quad Q=\operatorname{diag}(q_1,\ldots,q_n).
\]
Let $H=\vct{v}^{\perp}\subset\R^n$, let $P$ be the orthogonal projection
onto $H$, and let
\[
  \Lambda=P\Z^n,\qquad D=P[-1,1]^n,\qquad E=P\,Q^{1/2}B_2^n,
\]
where $B_2^n$ is the Euclidean unit ball.  $\Lambda$ is a lattice of rank
$d=n-1$ in $H$, with Euclidean covolume $\det\Lambda=1/\lVert\vct{v}\rVert_2$
(for primitive $\vct{v}$ the lattice $\Z^n\cap H$ has covolume
$\lVert\vct{v}\rVert_2$, and $P\Z^n$ is its dual in $H$).  $D$ is a zonotope
and $E$ is an origin-symmetric ellipsoid; we write $\Enorm{\cdot}$ for the
Euclidean norm on $H$ whose unit ball is $E$ and
$\langle\cdot,\cdot\rangle_E$ for its inner product.  Finally put
\begin{equation}\label{eq:F}
  F_n(\vct{x})^2=\prod_iq_i\cdot\sum_i\frac{x_i^2}{q_i},
  \qquad
  R_n(\vct{x})=\frac{n\,(\prod_ix_i)^{1/n}}{F_n(\vct{x})}.
\end{equation}

\subsection{The ellipsoid and the covolume}

\begin{lemma}[Inclusion]\label{lem:incl}
$E\subseteq D$.
\end{lemma}

\begin{proof}
For $\vct{y}\in H$ the support functions are $h_D(\vct{y})=\sum_i|y_i|$
(the support function of a projection is the restriction of the support
function) and $h_E(\vct{y})=\sup_{|\vct{b}|\leq1}\langle
\vct{y},PQ^{1/2}\vct{b}\rangle=|Q^{1/2}\vct{y}|=(\sum_iq_iy_i^2)^{1/2}$,
since $P\vct{y}=\vct{y}$.  It suffices to show $\sum_iq_iy_i^2\leq1$ on
the polytope $\{\vct{y}\in H:\sum|y_i|\leq1\}$, and since the left side
is convex, at its vertices.  These are the points where the edges of the
cross-polytope meet $H$; an edge from $e_i$ to $-e_j$ meets $H$ at
$(v_je_i-v_ie_j)/(v_i+v_j)$, and edges between two vectors of equal sign
do not meet $H$ because $\vct{v}$ is positive.  At such a vertex,
\[
  \sum_kq_ky_k^2=\frac{q_iv_j^2+q_jv_i^2}{(v_i+v_j)^2}\leq1
  \iff (q_i-1)v_j^2+(q_j-1)v_i^2\leq2v_iv_j .
\]
With $q_i-1=u_i(2-u_i)$ and $v_i=Mu_i$, $M=\max_jv_j$, dividing by
$M^2u_iu_j$ turns the right-hand inequality into
$u_j(2-u_i)+u_i(2-u_j)\leq2$, that is, $(1-u_i)(1-u_j)\geq0$.
\end{proof}

\begin{lemma}[Covolume in the ellipsoidal norm]\label{lem:det}
The covolume of $\Lambda$ with respect to $\langle\cdot,\cdot\rangle_E$ is
\[
  \covol_E(\Lambda)=\frac{1}{S\,F_n(\vct{x})}.
\]
\end{lemma}

\begin{proof}
Let $U$ be an $n\times(n-1)$ matrix whose columns form an orthonormal
basis of $H$ and put $\tilde{\vct{v}}=\vct{v}/\lVert\vct{v}\rVert_2$.  In
the coordinates $\vct{y}=U^{\mathsf T}\vct{h}$ on $H$, the polar body
$\{\vct{h}\in H:h_E(\vct{h})\leq1\}$ has Gram matrix $U^{\mathsf T}QU$ by
the computation of $h_E$ above, so $E$ itself has Gram matrix
$(U^{\mathsf T}QU)^{-1}$ and
$\covol_E(\Lambda)=\det\Lambda\cdot\det(U^{\mathsf T}QU)^{-1/2}$.  The
orthogonal matrix $O=[U\mid\tilde{\vct{v}}]$ gives
$\det Q=\det(O^{\mathsf T}QO)
 =\det(U^{\mathsf T}QU)\cdot\bigl(\tilde{\vct{v}}^{\mathsf T}Q\tilde{\vct{v}}
 -\tilde{\vct{v}}^{\mathsf T}QU(U^{\mathsf T}QU)^{-1}U^{\mathsf T}Q\tilde{\vct{v}}\bigr)$
by the Schur complement, and the bracket is the reciprocal of the
$(n,n)$ entry of $(O^{\mathsf T}QO)^{-1}=O^{\mathsf T}Q^{-1}O$, namely
$1/(\tilde{\vct{v}}^{\mathsf T}Q^{-1}\tilde{\vct{v}})$.  Hence
$\det(U^{\mathsf T}QU)=\det Q\cdot\tilde{\vct{v}}^{\mathsf T}Q^{-1}\tilde{\vct{v}}
=\prod_iq_i\cdot\sum_iv_i^2/q_i\big/\lVert\vct{v}\rVert_2^2$, and
\[
  \covol_E(\Lambda)^2
  =\frac{1}{\lVert\vct{v}\rVert_2^2}\cdot
   \frac{\lVert\vct{v}\rVert_2^2}{\prod_iq_i\sum_iv_i^2/q_i}
  =\frac{1}{S^2\prod_iq_i\sum_ix_i^2/q_i}
  =\frac{1}{S^2F_n(\vct{x})^2}.\qedhere
\]
\end{proof}

\subsection{Points with all coordinates far from integers}

The following lemma is Lemma~3.3 of Giri and Kravitz~\cite{GK26} in a
different normalization: in their notation it says
$\max\mathcal S_q(n)=\max\mathcal S_1(n+1-q)$, and $LRC(n+1-q)$ bounds
the right-hand side.  We include a short proof along their lines, since
the statement is used with every $q$ below.

\begin{lemma}[Separated points on rational subspaces]\label{lem:safe}
Let $1\leq q\leq n$ and assume $LRC(m)$ for every $m\leq n+1-q$.  Let
$V\subseteq\R^n$ be a rational subspace of dimension $q$ on which no
coordinate functional vanishes identically.  Then $V$ contains a point
$\vct{s}$ with
\[
  \nearest{s_i}\geq\frac{1}{n+2-q}\qquad(1\leq i\leq n).
\]
\end{lemma}

\begin{proof}
Induction on $q$, simultaneously for all $n$.  For $q=1$, $V=\R\vct{a}$
with $\vct{a}\in\Z^n$ primitive and all $a_i\neq0$.  The values $|a_i|$
are $m\leq n$ distinct positive integers;
$LRC(m)$ gives $t$ with $\nearest{ta_i}\geq1/(m+1)\geq1/(n+1)$, and
$\vct{s}=t\vct{a}$ works.

Let $q\geq2$.  The restrictions of the coordinate functionals to $V$ are
nonzero vectors $\vct{a}_1,\ldots,\vct{a}_n\in V$ (identify $V^{*}$ with
$V$ through the Euclidean inner product); they span $V^{*}$, so not all of
them are proportional.  Among all non-proportional pairs choose $(i,j)$
minimizing the angle between the lines $\R\vct{a}_i$ and $\R\vct{a}_j$,
and replace $\vct{a}_j$ by $-\vct{a}_j$ if necessary so that
$\langle\vct{a}_i,\vct{a}_j\rangle\geq0$; this amounts to changing the
sign of the $j$-th coordinate, which does not affect $\nearest{\cdot}$.
Put $\vct{h}=\vct{a}_i+\vct{a}_j\neq0$.  The line $\R\vct{h}$ makes a
strictly smaller angle with $\R\vct{a}_i$ than $\R\vct{a}_j$ does, so by
minimality $\vct{h}$ is proportional to no $\vct{a}_k$.  Let
$V'=\{\vct{y}\in V:\langle\vct{h},\vct{y}\rangle=0\}$, a rational
subspace of dimension $q-1$ on which $y_i=-y_j$ (after the sign change).
No coordinate vanishes on $V'$: the functional $\vct{a}_k$ vanishes on
$V'=\vct{h}^{\perp}\cap V$ only if $\vct{a}_k\in\R\vct{h}$.  Deleting the
$j$-th coordinate maps $V'$ injectively onto a rational subspace
$V''\subseteq\R^{n-1}$ of dimension $q-1$ with no vanishing coordinate.
The induction hypothesis (with $n-1$ in place of $n$; it requires
$LRC(m)$ for $m\leq(n-1)+1-(q-1)=n+1-q$) gives $\vct{s}''\in V''$ with all
coordinates at distance $\geq1/((n-1)+2-(q-1))=1/(n+2-q)$; its preimage
$\vct{s}\in V'$ has $\nearest{s_j}=\nearest{-s_i}=\nearest{s_i}$, so all
$n$ coordinates satisfy the bound.
\end{proof}

\subsection{Prefix inequalities for the Gram--Schmidt lengths}

Let $\vct{b}_1,\ldots,\vct{b}_d$ be any basis of $\Lambda$ ($d=n-1$), let
$\vct{b}_i^{*}$ be its Gram--Schmidt vectors with respect to
$\langle\cdot,\cdot\rangle_E$, and put $t_i=\Enorm{\vct{b}_i^{*}}^2>0$.
Then $\prod_it_i=\covol_E(\Lambda)^2$.

\begin{proposition}[Prefix inequalities]\label{prop:prefix}
Assume $LRC(m)$ for all $m<n$ and let $\vct{v}$ be a counterexample.  Then
for every $r=1,\ldots,d$,
\begin{equation}\label{eq:prefix}
  t_1+\cdots+t_r\;>\;B_r,\qquad
  B_r:=\Bigl(\frac{2r}{(n+1)(n+1-r)}\Bigr)^{2}.
\end{equation}
\end{proposition}

\begin{proof}
Fix $r$ and choose $\vct{z}_i\in\Z^n$ with $P\vct{z}_i=\vct{b}_i$
($1\leq i\leq r$).  The subspace $V=\spn(\vct{v},\vct{z}_1,\ldots,\vct{z}_r)$
is rational and has dimension $r+1$: if
$a\vct{v}+\sum a_i\vct{z}_i=0$, applying $P$ gives $\sum a_i\vct{b}_i=0$,
so all $a_i=0$ and then $a=0$.  No coordinate vanishes on $V$, since
$\vct{v}\in V$ has positive coordinates.  Lemma~\ref{lem:safe} with
$q=r+1$ (it needs $LRC(m)$ for $m\leq n-r\leq n-1$) gives
$\vct{s}=a\vct{v}+\sum_ia_i\vct{z}_i\in V$ with
$\nearest{s_k}\geq1/(n+1-r)$ for all $k$.

Apply nearest-plane rounding to the coefficients: choose integers
$k_r,k_{r-1},\ldots,k_1$ in this order so that
$\vct{e}:=\sum_{i\leq r}(a_i-k_i)\vct{b}_i=\sum_{i\leq r}\theta_i\vct{b}_i^{*}$
has $|\theta_i|\leq\frac12$ for all $i$ (at step $i$ the coefficient of
$\vct{b}_i^{*}$ is $a_i-k_i$ plus a quantity already fixed by
$k_{i+1},\ldots,k_r$, so $k_i$ can be chosen to bring it into
$[-\frac12,\frac12]$).  By orthogonality,
\begin{equation}\label{eq:round}
  \varepsilon:=\Enorm{\vct{e}}\leq\tfrac12\sqrt{t_1+\cdots+t_r}.
\end{equation}
By Lemma~\ref{lem:incl}, $\vct{e}\in\varepsilon E\subseteq\varepsilon D$,
so $\vct{e}=P\vct{y}$ for some $\vct{y}\in\R^n$ with $|y_k|\leq\varepsilon$
for all $k$.  The vector $\sum_i(a_i-k_i)\vct{z}_i-\vct{y}$ has projection
$\vct{e}-\vct{e}=0$, hence equals $c\vct{v}$ for some real $c$.  Therefore
\[
  \vct{s}-\sum_ik_i\vct{z}_i-\vct{y}=(a+c)\vct{v}=:t\vct{v},
\]
and since $\sum_ik_i\vct{z}_i\in\Z^n$,
\[
  \nearest{tv_k}=\nearest{s_k-y_k}\geq\nearest{s_k}-|y_k|
  \geq\frac{1}{n+1-r}-\varepsilon\qquad(1\leq k\leq n).
\]
If $\varepsilon\leq\frac1{n+1-r}-\frac1{n+1}=\frac{r}{(n+1)(n+1-r)}$,
then $t$ is a witness for $\vct{v}$, contradicting that $\vct{v}$ is a
counterexample.  So $\varepsilon>r/((n+1)(n+1-r))$, and
\eqref{eq:round} gives $t_1+\cdots+t_r\geq4\varepsilon^2>B_r$.
\end{proof}

For $r=1$ this is the finite-checking argument of Giri and
Kravitz~\cite[\S7]{GK26} in a sharper form: the geodesic
$t\mapsto t\vct{v}$ comes within $\varepsilon$ of every point of the
two-dimensional subtorus spanned by $\vct{v}$ and $\vct{z}_1$, in every
coordinate, and the $LRC(n-1)$ case---equivalently, the case of $n$
runners---applies to that subtorus.  For $r=d$ the inequality reads
$\sum_it_i>((n-1)/(n+1))^2$ and concerns the whole lattice.  The content
of \eqref{eq:prefix} is that an inequality of the same kind holds
simultaneously in every prefix of the flag, with a right-hand side that
grows with $r$.

\subsection{Korkine--Zolotarev bases and the product minimization}

\begin{lemma}[KZ bases]\label{lem:kz}
$\Lambda$ has a basis $\vct{b}_1,\ldots,\vct{b}_d$ whose Gram--Schmidt
lengths satisfy $t_{i+1}\geq\frac34t_i$ for $1\leq i<d$.
\end{lemma}

\begin{proof}
Construct the basis recursively (Korkine--Zolotarev~\cite{KZ1873}; see
\cite{LLS90}).  Let $\vct{b}_1$ be a shortest nonzero vector of $\Lambda$;
it is primitive.  Having chosen $\vct{b}_1,\ldots,\vct{b}_i$ spanning a
primitive sublattice, let $\pi_i$ be the orthogonal projection onto
$\{\vct{b}_1,\ldots,\vct{b}_i\}^{\perp}$, choose a shortest nonzero vector
of the projected lattice $\pi_i\Lambda$ (of rank $d-i$), and let
$\vct{b}_{i+1}\in\Lambda$ be any preimage; then $\vct{b}_{i+1}^{*}$ is that
shortest vector, and $\vct{b}_1,\ldots,\vct{b}_{i+1}$ again span a
primitive sublattice, so after $d$ steps we have a basis.  For the
inequality, note that $\vct{b}_i^{*}$ is a shortest vector of
$\pi_{i-1}\Lambda$, which also contains
$\pi_{i-1}(\vct{b}_{i+1}-m\vct{b}_i)=\vct{b}_{i+1}^{*}+(\mu-m)\vct{b}_i^{*}$
for every integer $m$, where $\mu$ is the Gram--Schmidt coefficient.
Choosing $m$ with $|\mu-m|\leq\frac12$ gives a nonzero vector of squared
length at most $t_{i+1}+\frac14t_i$, so $t_i\leq t_{i+1}+\frac14t_i$.
\end{proof}

Put $\rho=\frac34$, $B_0=0$, and $c_i=B_i-B_{i-1}$ for $1\leq i\leq d$.

\begin{lemma}[Product minimization]\label{lem:minimize}
Suppose $c_{i+1}/c_i>1/\rho$ for $1\leq i<d$.  If positive reals
$t_1,\ldots,t_d$ satisfy $t_{i+1}\geq\rho t_i$ for all $i<d$ and
$t_1+\cdots+t_r\geq B_r$ for all $r\leq d$, then
$\prod_it_i\geq\prod_ic_i$, with equality at $t_i=c_i$.
\end{lemma}

\begin{proof}
Put $w_i=\rho^{i-1}$ and $z_i=t_i/w_i$.  The constraints become
$z_1\leq z_2\leq\cdots\leq z_d$ and $\sum_{i\leq r}w_iz_i\geq B_r$
($1\leq r\leq d$); they define a polyhedron $\Pi\subset\R^d$ contained in
the open positive orthant ($z_1\geq B_1>0$).  The objective
$f(\vct{z})=\sum_i\log(w_iz_i)$ is concave on $\Pi$ and nondecreasing in
every coordinate, and the recession cone of $\Pi$,
$\{0\leq z_1\leq\cdots\leq z_d\}$, lies in the nonnegative orthant.
$\Pi$ contains no line, so every point of $\Pi$ is a convex combination of
vertices plus a recession vector, and $f$ is at least its minimum over the
vertices.  It remains to show $f\geq\sum\log c_i$ at every vertex.

At a vertex, $d$ linearly independent constraints are active.  The active
monotonicity constraints $z_i=z_{i+1}$ partition $\{1,\ldots,d\}$ into
maximal blocks of constant $z$; if there are $\beta$ blocks, then
$d-\beta$ monotonicity constraints are active and at least $\beta$ prefix
constraints must be active.  No prefix constraint is active at an index
$r$ strictly inside a block $\{a+1,\ldots,b\}$ ($a<r<b$) on which $z$
equals $\zeta$: feasibility at $a$ and activity at $r$ give
$\zeta\sum_{a<i\leq r}w_i\leq B_r-B_a=\sum_{a<i\leq r}c_i$, so $\zeta$
is at most a weighted average of the $c_i/w_i$ over $a<i\leq r$, hence
$\zeta\leq c_r/w_r$ because $c_i/w_i$ is strictly increasing; but
activity at $r$ and feasibility at $r+1$ give $w_{r+1}\zeta\geq c_{r+1}$,
i.e.\ $\zeta\geq c_{r+1}/w_{r+1}>c_r/w_r$, a contradiction.  Thus the
active prefixes are at block ends, there are at most $\beta$ of them, so
every block end is active, and on each block $\{a+1,\ldots,b\}$ the
$t_i=w_i\zeta$ form a geometric sequence with ratio $\rho$ whose sum is
$B_b-B_a=\sum_{a<i\leq b}c_i$.

Compare, within one block of length $L$, the decreasing geometric
sequence $\vct{g}=(t_{a+1},\ldots,t_b)$ with the sequence
$\vct{c}^{\downarrow}=(c_b,c_{b-1},\ldots,c_{a+1})$, both positive with
the same sum.  Consecutive ratios satisfy
$c^{\downarrow}_{j+1}/c^{\downarrow}_j<\rho=g_{j+1}/g_j$, so for $j\leq
k<l$ we have $c^{\downarrow}_l/c^{\downarrow}_j\leq g_l/g_j$ and
$c^{\downarrow}_j/c^{\downarrow}_k\geq g_j/g_k$, whence
\[
  \frac{\sum_{l>k}c^{\downarrow}_l}{\sum_{j\leq k}c^{\downarrow}_j}
  \leq\frac{\sum_{l>k}g_l/g_k}{\sum_{j\leq k}g_j/g_k}
  =\frac{\sum_{l>k}g_l}{\sum_{j\leq k}g_j}
  \qquad(1\leq k<L),
\]
and with equal totals, $\sum_{j\leq k}c^{\downarrow}_j\geq\sum_{j\leq
k}g_j$ for all $k$: $\vct{c}^{\downarrow}$ majorizes $\vct{g}$.  Since
$\sum\log$ is Schur-concave, $\prod_{a<i\leq b}c_i\leq\prod_{a<i\leq
b}t_i$.  Multiplying over the blocks gives $\prod_it_i\geq\prod_ic_i$ at
every vertex.  Finally $t_i=c_i$ is feasible ($c_{i+1}>c_i/\rho>\rho
c_i$ and all prefixes are equalities), so the bound is attained.
\end{proof}

For $n=14$ the minimum of the ratios $c_{i+1}/c_i$ is
$847/513>4/3$; for $n=15$ it is $3751/2349$, and for $n=13$ it is
$1175/688$.

\subsection{The shape factor}

\begin{lemma}[Shape bound]\label{lem:shape}
For $n=14$ and every $\vct{x}\in\R_{>0}^{14}$ with $\sum x_i=1$,
$R_{14}(\vct{x})<\frac12$.
\end{lemma}

\begin{proof}
Write $n=14$.  Passing from $x_i$ to $u_i=x_i/\max_jx_j$ cancels the
common scale in \eqref{eq:F}: with
\[
  q(u)=1+2u-u^2,\qquad a(u)=\frac{q(u)}{u^{2/n}},\qquad
  w(u)=\frac{u^2}{q(u)},\qquad
  H(\vct{u})=\Bigl(\sum_iw(u_i)\Bigr)\prod_ia(u_i),
\]
one has $R_n(\vct{x})^2=n^2/H(\vct{u})$, so it suffices to prove
$H>4n^2=784$ on the domain $\vct{u}\in(0,1]^n$ with some $u_i=1$.  Since
$w(1)=\frac12$, $w\geq0$, and $a(u)\to\infty$ as $u\to0$, $H$ tends to
infinity when any coordinate tends to $0$, so $H$ attains its minimum.

\emph{Interior coordinates.}  Fix all coordinates but one interior
coordinate $u\in(0,1)$ and put $C=\sum_{j\neq i}w(u_j)\in(0,\frac{n-1}2]$.
Up to a positive factor the objective is $\phi(u)=a(u)(C+w(u))$, and a
direct computation gives
\begin{equation}\label{eq:phi}
  \tfrac n2\,u^{1+2/n}\,\phi'(u)
  =-\bigl[C\{(n-1)u^2-(n-2)u+1\}-(n-1)u^2\bigr]=:-\beta(u).
\end{equation}
Let $\alpha$ be the smaller root of $(n-1)u^2-(n-2)u+1$ ($\alpha_{14}
=0.0926\ldots$).  Then $\beta(0)=C>0$ and $\beta(\alpha)=-(n-1)\alpha^2<0$.
If $C\leq1$, $\beta$ is concave with $\beta(0)>0$ and has exactly one
positive root; if $C>1$, $\beta$ is convex with $\beta(1)=2C-(n-1)\leq0$,
so its second root is $\geq1$.  In both cases $\beta$ has exactly one
root $u_0\in(0,\alpha)$ in $(0,1)$, and $\phi$ decreases on $(0,u_0)$ and
increases on $(u_0,1)$.  Hence at a minimizer every interior coordinate is
a critical point in $(0,\alpha)$.  At such a point, with
$A=C+w(u)=\sum_iw(u_i)$, the equation $\beta(u)=0$ is equivalent to
\begin{equation}\label{eq:crit}
  \frac{q(u)\{1-(n-2)u+(n-1)u^2\}}{n\,u^2(1+u)}=\frac1A .
\end{equation}
The derivative of the left side has the sign of
$-(13u^5+26u^4-50u^3-20u^2-7u+2)$, and the polynomial is positive on
$(0,\alpha)$ (drop its positive terms and use $\alpha<0.1$), so the left
side is strictly decreasing there.  Since $A$ is the same for every
coordinate, all interior coordinates of a minimizer share one value $t$.

\emph{Two or more coordinates equal to one.}  We use $a(u)>\frac85$ on
$(0,1]$: the logarithmic derivative of $a$ vanishes exactly at the roots
$(6\pm\sqrt{23})/13$ of $13u^2-12u+1$, so $a$ decreases from $+\infty$
to a local minimum at $u_1=(6-\sqrt{23})/13\in(0.092,0.093)$, increases
to a local maximum, and decreases to $a(1)=2$; hence
$a\geq\min(a(u_1),2)$ and $a(u_1)>q(0.092)/0.093^{1/7}>1.65$.  If $k\geq2$
coordinates equal one, then $\sum_iw(u_i)\geq k/2$ and
$H\geq\frac k2\,2^k(\tfrac85)^{n-k}$, which increases with $k$ and at
$k=2$ equals $4\cdot(8/5)^{12}>1125>784$.

\emph{Exactly one coordinate equal to one.}  Then the other thirteen equal
$t\in(0,\alpha)$ and
\[
  H(t)=\Bigl(\tfrac12+\tfrac{13t^2}{q(t)}\Bigr)\cdot2\cdot\frac{q(t)^{13}}{t^{26/14}}
  =\frac{(1+2t+25t^2)\,q(t)^{12}}{t^{13/7}} .
\]
On this interval its logarithmic derivative has the same sign as
$325t^3+23t^2+9t-1$, since all remaining factors are positive.  Thus $H$
decreases until the unique root
$t_0\in(0.0782,0.0783)$ of the increasing cubic and increases afterwards.
As $q$ and $1+2t+25t^2$ increase and $t^{-13/7}$ decreases,
\[
  H(t_0)\geq\frac{(1+2\cdot0.0782+25\cdot0.0782^2)\,q(0.0782)^{12}}{0.0783^{13/7}}
  >796.4>784 ,
\]
an exact rational comparison after raising to the seventh power.  So
$H>784$ in every case, and $R_{14}<14/28=\frac12$.
\end{proof}

The same argument gives $R_{15}<\frac12$ and $R_{13}<\frac{101}{200}$
(the corresponding cubics are $378t^3+25t^2+10t-1$ and
$276t^3+21t^2+8t-1$, and the constants $900$ and $(2600/101)^2$).  The
$n=14$ estimate is used in the proof, while the $n=15$ estimate is used
only in the outlook below.

\subsection{The theorem and its constants}

\begin{theorem}[Flag bound]\label{thm:flag}
Assume $LRC(m)$ for all $m<n$, let $\vct{v}\in\Z^n_{>0}$ be a primitive
counterexample to $LRC(n)$, and suppose $c_{i+1}/c_i>4/3$ for
$1\leq i<n-1$.  Put
\[
  K_n=\prod_{i=1}^{n-1}c_i,\qquad A_n=K_n^{-1/2}.
\]
Then $S\,F_n(\vct{x})\leq A_n$, and if moreover $R_n<\frac12$ on the
simplex, then
\[
  v_1\cdots v_n<\Bigl(\frac{A_n}{2n}\Bigr)^{n}.
\]
In particular, for $n=14$, using $LRC(13)$~\cite{Allikvere26} and
$LRC(m)$, $m\leq12$~\cite{ST26},
\[
  K_{14}=\frac{2289670297}{97243124257250844122250000000000000000},\qquad
  A_{14}=206083383792625.44\ldots,
\]
\[
  \log(v_1\cdots v_{14})<14\log\frac{A_{14}}{28}
  =414.77936455669981\ldots<\RawThreshUB .
\]
\end{theorem}

\begin{proof}
Take a KZ basis (Lemma~\ref{lem:kz}).  Its Gram--Schmidt lengths satisfy
the hypotheses of Lemma~\ref{lem:minimize} by Proposition~\ref{prop:prefix},
so $\covol_E(\Lambda)^2=\prod_it_i\geq K_n$, and Lemma~\ref{lem:det} gives
$(SF_n(\vct{x}))^{-2}\geq K_n$, i.e.\ $SF_n(\vct{x})\leq A_n$.  By the
definition \eqref{eq:F} of $R_n$,
\[
  n\,(v_1\cdots v_n)^{1/n}=n\,S\,(\textstyle\prod_ix_i)^{1/n}
  =S\,F_n(\vct{x})\,R_n(\vct{x})<\tfrac12A_n ,
\]
which is the product bound.  For $n=14$ the ratio hypothesis holds with
minimum $847/513$, the shape bound is Lemma~\ref{lem:shape}, and the
constants are the exact rational and algebraic numbers displayed.
\end{proof}

\begin{remark}[Comparison]\label{rem:compare}
The bound of Malikiosis--Santos--Schymura~\cite{MSS25} in the form used
in~\cite{ST26,Allikvere26} is $\sum_iv_i<\binom{n+1}{2}^{n-1}$, hence
$\log(v_1\cdots v_{14})<14(13\log105-\log14)=810.07\ldots$; the
reworking in~\cite{BCS26} gives the same bound, and the bound
of~\cite{GK26} is weaker.  Theorem~\ref{thm:flag} lowers the threshold
to $414.78$.  Three ingredients account for the gain.  Using all $n-1$
prefix inequalities of a flag, rather than a single density statement, is
the largest step: with the successive minima of the projected cube in
place of $t_i$ and Minkowski's second theorem in place of
Lemma~\ref{lem:det}, the same prefix inequalities already give
$\sum_iv_i<(n+1)!\,n!/2^n$ and a threshold of $570.46$
(Table~\ref{tab:thresholds}).  Passing from successive minima to the
Gram--Schmidt lengths of a KZ basis makes the covolume exact and removes
Minkowski's factor $2^{n-1}$ together with any Hermite-type constant.
The adapted ellipsoid $E$ removes the $\sqrt n$-type loss of converting a
Euclidean estimate back into the sum of the speeds; its shape factor is
controlled by Lemma~\ref{lem:shape}.  For $n=13$ the theorem gives
$13(\log A_{13}-\log13+\log\frac{101}{200})=341.03\ldots$ against
$670.35$, so the first $59$ verified primes in increasing order
of~\cite{Allikvere26} (through $p=523$, with logarithmic sum
$341.17\ldots$) already suffice for fourteen runners.
\end{remark}

\begin{table}[ht]\centering
\caption{Logarithmic finite-checking thresholds for $\log(v_1\cdots v_n)$,
assuming $LRC(m)$ for $m<n$.  Column two is the bound of~\cite{MSS25} in
the form $B_n=(\binom{n+1}{2}^{n-1}/n)^n$ used in~\cite{ST26,Allikvere26};
column three uses the prefix inequalities~\eqref{eq:prefix} with
successive minima and Minkowski's second theorem; column four is
Theorem~\ref{thm:flag}.  The forced divisor of Lemma~\ref{lem:divisor} is
not included.}
\label{tab:thresholds}
\begin{tabular}{@{}cccc@{}}\toprule
$n$ & $\log B_n$~\cite{MSS25} & successive minima & Theorem~\ref{thm:flag}\\
\midrule
$13$ & $670.35$ & $470.18$ & $341.03$\\
$14$ & $810.07$ & $570.46$ & $414.78$\\
$15$ & $964.75$ & $681.99$ & $497.03$\\
\bottomrule
\end{tabular}
\end{table}

\subsection{The forced divisor}

\begin{lemma}[Forced divisor]\label{lem:divisor}
If $\vct{v}\in\Z_{>0}^{n}$ is a counterexample to $LRC(n)$, then the
integer $\lcm(2,3,\ldots,n+1)$ divides $v_1\cdots v_n$.  For $n=14$,
\[
  \lcm(2,\ldots,15)=360360=2^3\cdot3^2\cdot5\cdot7\cdot11\cdot13,\qquad
  \log360360=12.79485\ldots>\DivisorLogLB .
\]
\end{lemma}

\begin{proof}
Fix $2\leq d\leq n+1$.  If no $v_i$ is divisible by $d$, then at
$t=1/d$ every $\nearest{v_i/d}$ is a nonzero multiple of $1/d$, hence
$\geq1/d\geq1/(n+1)$, and $t$ is a witness.  So every $d\leq n+1$
divides some $v_i$.  For each prime $r\leq n+1$, let $q_r$ be the largest
power of $r$ not exceeding $n+1$.  Then $q_r$ divides the speed product.
The numbers $q_r$ are pairwise coprime, and their product is
$\lcm(2,\ldots,n+1)$.
\end{proof}

\begin{corollary}[Sufficient logarithmic prime sum]\label{cor:mass}
Let $\mathcal P$ be a set of primes $p>13$ such that, for every $p\in
\mathcal P$, every primitive counterexample
$\vct{v}\in\Z_{>0}^{14}$ to $LRC(14)$ with distinct coordinates has
$p\mid v_1\cdots v_{14}$.  If
$\sum_{p\in\mathcal P}\log p>\TargetUB$, then $LRC(14)$ holds.
\end{corollary}

\begin{proof}
Let $\vct{v}$ be a counterexample.  As in~\cite{Allikvere26}, we may take
the coordinates positive, pairwise distinct (if two coincide, at most
thirteen distinct speeds occur and $LRC(13)$ gives a witness with
$\nearest{tv_i}\geq\frac1{14}>\frac1{15}$), and primitive (scaling
preserves the lonely-runner property).  Then $360360\prod_{p\in\mathcal
P}p$ divides $v_1\cdots v_{14}$, the factors being pairwise coprime, so
\[
  \log(v_1\cdots v_{14})\geq\log360360+\sum_{p\in\mathcal P}\log p
  >\DivisorLogLB+\TargetUB=\RawThreshUB>14\log\frac{A_{14}}{28},
\]
contradicting Theorem~\ref{thm:flag}.
\end{proof}

\section{Verification for one prime}\label{sec:pipeline}

Fix a prime $p$ and set $k=14$.  The verification has four stages.  We
generate $I(14,p,1)$ up to units, apply successive binary lifts, examine
the surviving orbits exactly at level $15$, and record whether the
prime-divisibility criterion holds.

We call the elements of $\Z_p^{\times}/\{\pm1\}$ \emph{classes} and
represent them by $1,\ldots,(p-1)/2$.  A \emph{row} is a $14$-multiset of
classes representing an element of $I(14,p,1)$ up to permutation and sign
changes.  A speed class $v$ \emph{covers} a time class $a$ when
$15\,d_p(av)<p$, where $d_p(x)=\min(r,p-r)$ for
$x\equiv r\pmod p$ and $0\leq r<p$.

The binary-lifting stages compute $F_l(r)$ for $l=2,4,\ldots,32$.  We say
that a unit orbit \emph{survives the binary lifts} if
$F_{32}(r)\neq\varnothing$.  The final stage determines $F_{15}(r)$
exactly for each surviving orbit.

\subsection{The \texorpdfstring{level-$1$}{level-1} family}\label{sec:basefamily}

At level one a row is improper exactly when its fourteen classes jointly
cover every time class.  Every improper row either consists of fourteen
distinct classes none of which can be removed (\emph{irredundant
branch}), or contains a cover on at most thirteen distinct classes
(\emph{reducible branch}).  Let $n=(p-1)/2$, let
$m=\lfloor(p-1)/15\rfloor$ be the number of time classes covered by one
speed class, and let $\tau(p)$ be the smallest number of speed classes
that cover all time classes.

\begin{lemma}[Private times]\label{lem:quota}
In any cover of the $n$ time classes by at most $s$ speed classes, some
speed class covers at least $q_s=\lceil\max(0,2n-sm)/s\rceil$ time
classes that no other class of the cover covers (\emph{private} times).
If the cover is the support of a row on at most $13$ distinct classes,
such a class can be chosen with multiplicity at most two in the row, and
$q_{13}\geq1$.
\end{lemma}

\begin{proof}
As in Lemma~3.1 of~\cite{Allikvere26}: if $n_j$ time classes are covered
exactly $j$ times and the cover has $s'\leq s$ classes, then
$s'm\geq n_1+2(n-n_1)$, so $n_1\geq2n-sm$ private times are shared among
$s'$ classes.  For the second statement let $h$ be the number of support
classes of multiplicity $\geq3$; then $3h+(s'-h)\leq14$, so
$s'+h\leq\lfloor(14+s')/2\rfloor\leq13$ for $s'\leq13$.  Since
$13m\leq\frac{13}{15}(p-1)<2n$, we have $2n-13m>0$; the $h$ heavy classes
own at most $hm$ private times, so the classes of multiplicity $\leq2$ own
at least $2n-s'm-hm\geq2n-13m>0$ of them, and one of these $s'-h\leq13$
classes owns at least $q_{13}$.
\end{proof}

\paragraph{Decomposition variant ($\tau(p)\geq13$).}
If no cover on $\leq12$ classes exists, every $13$-class cover is minimal,
and the reducible branch is enumerated as the irredundant $13$-class
covers containing the class $1$ (mode \texttt{km1root}), each extended by
one arbitrary class.  The precondition $\tau(p)\geq13$ is verified before
generation by an exhaustive search for covers on $\leq12$ classes (mode
\texttt{km1low 12}), pruned by Lemma~\ref{lem:quota}: in a cover on
$s\leq12$ classes some class has $q_{12}$ private times, and a unit moves
it to class~$1$.  The node count and the verdict of this search are part
of the certificate.

\paragraph{General variant ($\tau(p)\leq12$).}
If \texttt{km1low 12} finds covers, non-minimal $13$-class covers exist
and the decomposition is incomplete.  The generator then runs the mode
\texttt{km1all}: all supports on $\leq13$ classes containing class $1$
that pass the quota $q_{13}$, with class $1$ of multiplicity one, all
multiplicity vectors, and all single-class extensions.

\begin{lemma}[Completeness of the general variant]\label{lem:km1all}
Every improper row $M$ is produced by the irredundant branch or by
\texttt{km1all}, up to the action of units.
\end{lemma}

\begin{proof}
If $M$ has fourteen distinct classes none of which is removable, it lies
in the irredundant branch.  If $M$ has support on $s'\leq13$ classes, let
$c$ be the class of multiplicity $\leq2$ with $\geq q_{13}$ private times
given by Lemma~\ref{lem:quota}; remove one copy of $c$ if its
multiplicity is two, otherwise one copy of any repeated class.  The
result is a $13$-multiset $C$ with the same support in which $c$ has
multiplicity one and $\geq q_{13}$ private times; a unit moves $c$ to
class $1$, so $C$ is enumerated, and one extension restores $M$.  If $M$
has fourteen distinct classes and a removable class, remove it; the
remaining $13$-class cover has a class with $\geq q_{13}$ private times by
the first part of Lemma~\ref{lem:quota} and all multiplicities one, so it
is enumerated after that class is normalized to $1$, and adding the
removed class recovers $M$.
\end{proof}

The inclusion \texttt{km1low}$\subseteq$\texttt{km1all}, checked by the
verification script, is a consistency test rather than the completeness
proof.  The precondition fails at the \NearGcdCount{} verified primes with
$\tau(p)\leq12$ in Table~\ref{tab:gates}, where the general variant is used.  It
holds at the other \NoNearGcdCount{} primes, where the decomposition
variant is used.  Appendix~\ref{app:audit} records the corresponding
certificate labels and additional covering-number checks.

\paragraph{Generation up to units.}
The irredundant branch emits one representative per orbit of
$\Z_p^{\times}$, the lexicographic minimum among the normalizations to
class $1$ that pass the quota; by Lemma~\ref{lem:quota} every orbit has
such a normalization, and by Proposition~5.1 of~\cite{ST26} one
representative suffices.  The reducible branch may emit several
representatives of one orbit; before the level-$15$ step, rows that survive
the binary lifts are reduced to the lexicographically least element of
their unit orbit
(in parallel, since there are millions of rows when $\tau(p)\leq12$).

\subsection{Successive binary lifting}

Each row has $2^{14}$ lifts from level $1$ to level $2$; the filter tests
every lift exactly (witnesses in $(2p)^{-1}\Z$ by bit masks, and the gcd
condition), retains the improper ones, and repeats with $c=2$ through
levels $4,8,16,32$.  By the invariant of Section~\ref{sec:framework} the
set carried at level $l$ is $F_l(r)$, so the row is discarded when its
improper fiber $F_l(r)$ is empty.  At every verified prime with
$\tau(p)\geq13$, the only orbit that survives through level $32$ is
represented by
\[
  \vct{a}=(1,2,\ldots,14)
\]
and every other row has empty improper fiber by level $8$.  When
$\tau(p)\leq12$, between $8{,}011$ and $8{,}235{,}157$ unit orbits survive
through level $32$, so binary lifting alone is not enough.  The detailed
histograms and exceptional rows are recorded in Appendix~\ref{app:audit}.

\subsection{\texorpdfstring{Level-$15$ elimination by factoring
$15=3\cdot5$}{Level-15 elimination by factoring 15=3*5}}

By Remark~3.2 of~\cite{ST26} a level divisible by $k+1=15$ is unavoidable
for $\vct{a}$.  We therefore examine every orbit representative surviving
the binary lifts at level $15$, where its fiber consists of the $15^{14}$ lifts
$w_i=v_i+a_ip$, $a_i\in\{0,\ldots,14\}$.  Coordinatewise branch-and-bound,
the method of~\cite{Allikvere26} for $7^{13}$ lifts, would face a fiber
$3.0\times10^5$ times larger and would have to be repeated for the
$10^4$--$10^7$
orbits arising when $\tau(p)\leq12$.  Instead, we exploit the
factorisation $15=3\cdot5$.

\begin{lemma}[Witnesses lift along divisors of the level]\label{lem:factor}
Let $f\mid15$, let $\vct{w}$ be a lift of $r$ to level $15$, and let
$\vct{w}'\equiv\vct{w}\pmod{fp}$ be its reduction to level $f$.  If
$\vct{w}'$ has a witness $t\in(fp)^{-1}\Z$, then $t$ is a witness for
$\vct{w}$.  Consequently every witness-free level-$15$ lift of $r$
reduces to a witness-free level-$3$ lift and a witness-free level-$5$
lift, and is determined by them through the Chinese remainder theorem on
the digits $a_i$.
\end{lemma}

\begin{proof}
$t(w_i-w_i')\in t\cdot fp\Z\subseteq\Z$, so
$\nearest{tw_i}=\nearest{tw_i'}\geq\frac1{15}$, and
$(fp)^{-1}\Z\subseteq(15p)^{-1}\Z$.  The digits satisfy
$a_i\equiv a_i'\pmod f$, and $(a_i\bmod3,a_i\bmod5)$ determines $a_i$.
\end{proof}

The witness-free lifts at levels $3$ and $5$ are found by depth-first
searches over $3^{14}$ and $5^{14}$ digit vectors, branching on the time
class with fewest remaining options; each Chinese-remainder candidate is
then tested at level $15$ against the full grid $(15p)^{-1}\Z$ and the
gcd condition $\gcd(15,w_j:j\neq i)>1$ for some $i$.  Thus $F_{15}(r)$ is determined
exactly.  For $\vct{a}$, the level-$3$ and level-$5$ searches leave three
and five witness-free lifts, respectively.  They combine to fifteen
Chinese-remainder candidates, seven of which remain witness-free at level
$15$; all seven satisfy the gcd condition.  Four have every coordinate
divisible by $3$, two have every coordinate divisible by $5$, and the last
is $15$ times a unit multiple of $\vct{a}$.  This calculation takes about
one second.  Appendix~\ref{app:audit} records the independent checks of the
final stage.

\subsection{The few-exception shift lemma when \texorpdfstring{$\tau(p)\leq12$}{tau(p) <= 12}}

At the primes with $\tau(p)\leq12$, successive binary lifting still leaves
between $8{,}011$ and $8{,}235{,}157$ surviving unit orbits.  A typical
orbit comes from a $12$-class cover.  Its witness-free level-$15$ lifts
have twelve coordinates divisible by $15$ and two unrestricted
coordinates; the gcd condition can disregard only one coordinate and
therefore need not make such a lift proper.  The following lemma supplies
the remaining step in the direct divisibility argument.

\begin{lemma}[Few-exception shift]\label{lem:neargcd}
Let $u_1,\ldots,u_{14}$ be positive integers, let $d$ be a prime, let
$e=\#\{i:d\nmid u_i\}$, and put $m_d=\lceil2d/15\rceil$.  Assume
$LRC(m)$ for all $m\leq12$.  If $2\leq e$ and $e\,m_d<d$, then there is a
rational $t$ with $\nearest{tu_i}\geq\frac1{15}$ for all $i$.  In
particular, for a vector that does not satisfy the level-$15$ gcd condition,
these hypotheses hold when at least twelve coordinates are divisible by
$3$, or at least ten are divisible by $5$.
\end{lemma}

\begin{proof}
Let $B$ be the block of the $14-e\leq12$ speeds divisible by $d$.  After
repetitions are removed, these speeds take $m\leq14-e\leq12$ distinct
values.  The statement $LRC(m)$ gives $t_0$ with
$\nearest{t_0u_i}\geq\frac1{13}>\frac1{15}$ on $B$.  Consider
$t_q=t_0+q/d$, $q=0,\ldots,d-1$.  Every $u_i\in B$ is fixed modulo one by
these shifts.  For an exceptional $u$, the positions $t_0u+qu/d$
($q=0,\ldots,d-1$) form a full $d$-point grid because $d$ is prime and
$d\nmid u$; an open arc of length $2/15$ contains at most
$\lceil2d/15\rceil=m_d$ points of a $d$-grid.  So each exception forbids
at most $m_d$ shifts, all $e$ exceptions forbid at most $em_d<d$, and some
$t_q$ is good for every speed.  The set of good $t_0$ is a finite union of
intervals with rational endpoints, so $t_0$, hence $t_q$, may be taken
rational.  In the two cases stated in the last sentence of the lemma,
$m_3=m_5=1$, with $e\leq2$ and $e\leq4$, respectively.
\end{proof}

The level-$15$ program applies Lemma~\ref{lem:neargcd} with $d=3$ and
$d=5$ to every witness-free lift left improper by the gcd condition.
Divisibility by $3$ and $5$ is constant on a residue class modulo $15p$,
so the conclusion holds for every integer vector in that class.  This is
the last alternative in Lemma~\ref{lem:gate}(ii).  It proves the required
divisibility by $p$, but not necessarily $J(14,p)=\varnothing$.  The input
$LRC(m)$ for $m\leq12$ is Theorem~1.3 of~\cite{ST26}.

\begin{proposition}[Verified prime-divisibility computations]\label{prop:closure}
For every prime $p$ in Table~\ref{tab:gates}, every primitive
counterexample $\vct{v}\in\Z_{>0}^{14}$ to $LRC(14)$ with distinct
coordinates has $p\mid v_1\cdots v_{14}$.  For the primes in the second
block of Table~\ref{tab:gates}, moreover $J(14,p)=\varnothing$.
\end{proposition}

\begin{proof}
Every element of $I(14,p,1)$ lies in the unit orbit of a generated row
(Lemmas~\ref{lem:quota} and~\ref{lem:km1all}, with the precondition
verified where the decomposition variant was used).  A row with
$F_l(r)=\varnothing$ at some binary level is eventually proper.  For a
row that survives the binary lifts, Lemma~\ref{lem:factor} shows that the
level-$15$ search finds every witness-free lift.  If all such lifts satisfy
the gcd condition, then $F_{15}(r)=\varnothing$, and
Lemma~\ref{lem:gate}(i) applies; equivalently,
$I(14,p,480)=\varnothing$ at the common level $\lcm(32,15)$.  If some
residue classes are handled only by Lemma~\ref{lem:neargcd}, they satisfy
the third alternative of Lemma~\ref{lem:gate}(ii) at level $15$.  In both
cases Lemma~\ref{lem:gate} gives $p\mid v_1\cdots v_{14}$.
\end{proof}

Appendix~\ref{app:examples} gives one complete example of each generator
variant, including the recorded counts and running times.

\section{The verified primes and the proof of \texorpdfstring{$LRC(14)$}{LRC(14)}}\label{sec:gateset}

Let $P$ be the set of \GateCount{} verified primes listed in
Table~\ref{tab:gates}.

\begin{longtable}{@{}>{\raggedright\arraybackslash}p{0.20\textwidth}>{\raggedright\arraybackslash}p{0.50\textwidth}>{\raggedleft\arraybackslash}p{0.08\textwidth}>{\raggedleft\arraybackslash}p{0.12\textwidth}@{}}
\caption{Verified primes.  Logarithms are natural, and block sums are rounded down to four decimals.  The first block has $\tau(p)\leq12$ and uses the direct divisibility form of Lemma~\ref{lem:gate}, with Lemma~\ref{lem:neargcd}; the second block verifies $J(14,p)=\varnothing$.}\label{tab:gates}\\
\toprule Block & Primes & Count & $\sum\log p$\\ \midrule \endfirsthead
\multicolumn{4}{l}{\small\itshape Table~\ref{tab:gates}, continued}\\ \toprule Block & Primes & Count & $\sum\log p$\\ \midrule \endhead
\midrule \multicolumn{4}{r}{\small\itshape Continued on next page}\\ \endfoot \bottomrule \endlastfoot
Low covering number, $\tau(p)\leq12$ & 89, 131, 149, 157, 163, 167, 173, 179, 181, 191, 193, 197, 199, 211, 223, 227, 229, 233, 241 & 19 & 98.8199\\[2pt]
$\tau(p)\geq13$, $J(14,p)=\varnothing$ & 239, 251, 257, 263, 269, 271, 277, 281, 283, 293, 307, 311, 313, 317, 331, 337, 347, 349, 353, 359, 367, 373, 379, 383, 389, 397, 401, 409, 419, 421, 431, 433, 439, 443, 449, 457, 461, 463, 467, 479, 487, 491, 499, 503, 509, 521, 523, 541, 547, 557, 563, 569 & 52 & 310.0033\\[2pt]
\textbf{Total} & & \textbf{71} & \textbf{408.8233}\\
\end{longtable}


Evaluated at fifty-digit precision,
\begin{equation}\label{eq:mass}
  \sum_{p\in P}\log p>\GateMassLB>\TargetUB .
\end{equation}
The margin above the target is \MassMargin{},
more than $\log p$ for any single $p\in P$.

\subsection{Proof of the main theorem}

\begin{proof}[Proof of Theorem~\ref{thm:main}]
The cases $m\leq7$ are known, with $LRC(7)$ proved by
Rosenfeld~\cite{Rosenfeld25}; $LRC(8)$ and $LRC(9)$ are proved
in~\cite{Trakulthongchai26}.\footnote{The chain passes through $LRC(9)$.
The verification program published with~\cite{Trakulthongchai26} forms
the intersection of two lifted families with a key that depends on the
coordinate order, so the set it passes to the final check may be a proper
subset of the full intersection; the author has been informed, and a
corrected version is forthcoming.  The emptiness of $J(9,p)$ is
established independently, without an intersection step, by the
successive-lifting code accompanying~\cite{ST26};
Appendix~\ref{app:source} records our check of that computation.}  The
cases $10\leq m\leq12$ follow from Theorem~1.3 of~\cite{ST26}, and
$LRC(13)$ from~\cite{Allikvere26}.  By
Proposition~\ref{prop:closure} and~\eqref{eq:mass}, the hypotheses of
Corollary~\ref{cor:mass} are satisfied.  Hence $LRC(14)$ holds.
\end{proof}

\section{Computation and reproducibility}

Each verified prime has a certificate recording the generator variant,
the covering-search counts, every binary-lifting stage, and the exact
level-$15$ calculation.  A separate script checks these records against
the claimed conclusion, verifies the retained files by SHA-256, and
recomputes the logarithmic prime sum.  Two further scripts derive the
analytic constants from their definitions.  The code, certificates, and
verification reports are archived at the DOI given below.  The appendices
give one complete example of each generator variant, the detailed checks
and their limits, and the computational cost.

\section{Discussion}\label{sec:universality}

\subsection{Orbits surviving the binary lifts}

A speed tuple of length $k$ is \emph{tight} if
$\max_{t\in\R}\min_i\nearest{t v_i}=1/(k+1)$.

At each of the \NoNearGcdCount{} verified primes with $\tau(p)\geq13$, the
only unit orbit to survive the binary lifts is that of
$\vct{a}=(1,\ldots,14)$.  Its seven witness-free level-$15$ lifts are all
proper by the gcd condition.  In the $k=13$ computation, two orbits survived
at every verified prime: $(1,\ldots,13)$ and the Goddyn--Wong tight set
$(1,\ldots,11,13,24)$~\cite{GW06,Allikvere26}.

For the verified $k=14$ primes with $\tau(p)\geq13$, every other row has
$F_8(r)=\varnothing$.  This includes
$\vct{c}=(1,\ldots,11,13,14,24)$.  At $t=12/29$, every coordinate of
$\vct{c}$ is at distance at least $2/29$ from an integer.  Since
$2/29>1/15$, this tuple is not tight.  The survival of $\vct{a}$ through
the binary levels is forced by Remark~3.2 of~\cite{ST26}: modulo one, its
witness times are $s/15$ with $\gcd(s,15)=1$, whereas none of the levels
$2,4,8,16,32$ is divisible by $15$.  The computational observation is
that no other orbit survives at these primes.

This resembles Proposition~7.1 of~\cite{ST26}, although that result
concerns level-$1$ witnesses for all sufficiently large primes; we make no
asymptotic claim here.  When $\tau(p)\leq12$, the family surviving the
binary lifts is much larger and is handled by the direct divisibility
argument of Lemma~\ref{lem:neargcd}, rather than by emptying an improper
fiber at a grid level.

\subsection{Outlook: sixteen runners}\label{sec:outlook}

Theorem~\ref{thm:flag} applies to $n=15$ now that $LRC(14)$ is available.
The ratio hypothesis holds with minimum $3751/2349$, and the proof of the
shape estimate also gives $R_{15}<\frac12$.  Thus
$15\log(A_{15}/30)=497.03\ldots$.  Since
$\log\lcm(2,\ldots,16)=\log720720$, the logarithmic sum required from
verified primes is
\[
  497.03\ldots-\log720720=483.54\ldots.
\]
The verification programs are
generic in $K$, but the terminal level $16=2^4$ is itself a binary level,
so the factorisation used for level $15$ does not apply directly.  We
record these bounds as a starting point; the computation is future work.

\section{Conclusion}

The full flag of prefix inequalities gives a substantially smaller
finite-checking bound.  For fourteen speeds the logarithmic threshold
falls from $810.07$ to $414.78$, and the forced divisor reduces the
required logarithmic prime sum further to $\TargetUB$.  The
\GateCount{} verified primes exceed this threshold and prove $LRC(14)$,
equivalently the Lonely Runner Conjecture for fifteen runners.

\section*{Data availability}

The certificates, evidence packages, source code, checker scripts, and
verification reports are archived on Zenodo,
\url{https://doi.org/10.5281/zenodo.22667683}.

\section*{Use of artificial intelligence}

The author used large language models for code development, drafting,
proof exploration, and internal adversarial readings.  No model output is
mathematical evidence:
proofs are given in full, and exact scripts and archived certificates
support the numerical and computational claims.  The author checked the
arguments, code, and data and takes full responsibility.

\appendix

\section{Two complete verification examples}\label{app:examples}

\paragraph{$p=547$ (decomposition variant).}
The precondition search \texttt{km1low 12} visited $453{,}242{,}538$ nodes
in $774$ seconds and found no cover on at most $12$ classes.  The
irredundant branch ran $1{,}032$ jobs, visited
$1{,}425{,}111{,}896{,}802$ nodes, and produced two rows; the $36$
\texttt{km1root} jobs found no $13$-class cover.  During binary lifting,
the improper fiber of $(1,\ldots,11,13,14,24)$ is empty at level $4$,
whereas $\vct{a}$ survives through level $32$.  The level-$3$ and level-$5$
searches for $\vct{a}$ produce three and five witness-free lifts,
respectively, hence $15$ Chinese-remainder candidates.  Seven remain
witness-free at level $15$, and all seven are proper by the gcd condition.
Generation used $453{,}624$ CPU-seconds, or $126.0$ CPU-hours, and took
$7.9$ hours on one sixteen-vCPU machine.  Of the $1{,}032$ jobs, $297$
were repeated on other hosts with identical rows and node counts.  Thus
the computation verifies $J(14,547)=\varnothing$.

\paragraph{$p=233$ (general variant).}
Here $\tau(233)=12$: \texttt{km1low 12} found $24$ covers on $12$ classes
up to units.  The irredundant branch ran $178$ jobs, visited
$7{,}408{,}758{,}267$ nodes, and produced $5{,}767{,}782$ rows in $1{,}861$
CPU-seconds.  The \texttt{km1all} search found $6{,}746$ covers in $238$
CPU-seconds and produced $782{,}536$ extended rows.  Of the $85{,}616$ rows
with $F_2\neq\varnothing$, $58{,}229$ have $F_4=\varnothing$ and $81$ more
have $F_8=\varnothing$.  The remaining $27{,}306$ rows represent
$13{,}457$ unit orbits that survive through level $32$.

At level $15$, these orbits have $3{,}674{,}453$ witness-free lifts.
The gcd condition leaves $1{,}056{,}768$ of them improper, and every one
of the corresponding integer residue classes satisfies
Lemma~\ref{lem:neargcd}.  The orbit of $\vct{a}$ is handled by the gcd
condition alone; the other $13{,}456$ orbits use the lemma.  The
level-$15$ stage used $881$ CPU-seconds, and the complete verification took
about $0.85$ CPU-hours and ten minutes on one machine.  Thus the computation
verifies the divisibility conclusion $233\mid v_1\cdots v_{14}$ required
by Lemma~\ref{lem:gate}(ii).

\section{Detailed verification record}\label{sec:audit}\label{app:audit}

\subsection{Certificates}
Each verified prime has an evidence package \texttt{evidence\_P\_s0.tgz}.
Its \texttt{SUMMARY.json} records the prime, generator variant,
precondition result and node count, the job, node, and row counts for both
branches, the binary-lifting histogram, and the surviving orbit
representatives with their multiplicities.  For each orbit it also records
the level-$3$, level-$5$, and Chinese-remainder counts, the number of
witness-free level-$15$ lifts, and the two final counters
\texttt{improper\_after\_gcd} and
\texttt{improper\_after\_neargcd}.  The second name is implementation
shorthand: it counts residue classes not handled by
Lemma~\ref{lem:neargcd}, rather than membership in another modular
improper set.

The same package contains a SHA-256 manifest, the precondition output,
per-job generator output and statistics, the row files from both branches,
the binary-lifting output, every level-$15$ log, and the run log
\texttt{run\_P.log}.  A \texttt{SHA256SUMS.txt} covers the packages made on
each machine.  The archive also includes the source of the generator,
binary-lifting program, level-$15$ program, driver, and checker for the
analytic constants.  Every manifest records the SHA-256 hashes of the
three executables used for the computation; these hashes are the same for
all \GateCount{} primes.  A fresh build on a machine of the same type
produced byte-identical executables.

The Python standard-library script \texttt{audit\_lrc14\_gates.py} checks
each certificate against its claim, re-hashes its manifest, and recomputes
the prime sum at sixty-digit precision.  The decomposition precondition
fails for the \NearGcdCount{} primes through $241$ in the first block of
Table~\ref{tab:gates}; a separate search gives $\tau(p)=12$ for every
prime $139\leq p\leq233$.  The variant labels are \texttt{general} and
\texttt{decomposition}; the final claim labels are \texttt{divisibility}
and \texttt{J(K,p) empty}.

\subsection{Binary-lifting summary}

For the primes with $\tau(p)\geq13$, the number of generated rows with
$F_2(r)\neq\varnothing$ ranges from one (for most primes above $480$) to
$26{,}773$ at $p=271$.  Every such row except $\vct{a}$ has
$F_4(r)=\varnothing$, apart from one row at each of
$p=239,271,311,317,331,337,349,421$; these eight rows have
$F_8(r)=\varnothing$.  No row other than $\vct{a}$ reaches level $16$.
The row $(1,\ldots,11,13,14,24)$ is the level-$8$ row at $p=239$ and the
only other row with $F_2(r)\neq\varnothing$ at $p=547$, where
$F_4(r)=\varnothing$.

For the primes with $\tau(p)\leq12$, the number of unit orbits surviving
through level $32$ ranges from $8{,}011$ at $p=179$ to $8{,}235{,}157$ at
$p=181$.  A few additional rows reach level $16$ before their improper
fibers become empty.

\subsection{Checks}
\begin{enumerate}[leftmargin=*,label=(\roman*)]
\item \emph{Analytic constants.}  \texttt{kz\_flag\_check.py} recomputes
  $B_r$, $c_i$, the ratios $c_{i+1}/c_i$, $K_{14}$, $A_{14}$, the
  threshold, and the divisor credit in exact rational arithmetic and at
  fifty digits; an optional mode enumerates the $4096$ vertices of the
  polyhedron of Lemma~\ref{lem:minimize} and confirms the minimum; a
  second script re-derives the same numbers from the definitions.  Every
  strict inequality between displayed decimal values in the article is
  checked at fifty-digit precision.  The
  polynomial identities in the proof of Lemma~\ref{lem:shape} were checked
  symbolically.  As a numerical consistency check, four hundred random
  starts found no value of $H$ below $802.5$, while the proof gives the
  lower bound $784$.
\item \emph{Generator.}  The engine is the generator of~\cite{Allikvere26}
  compiled at $K=14$, with a faster search kernel and an additional mode
  for the reducible branch.  Against the build without the new kernel it produces
  identical rows and node counts at $p=83$ (full census, $101{,}974$ rows)
  and $p=431$ (all $616$ irredundant jobs).
  At $p=239$, a separately written probe reproduces the recorded level-$1$
  and level-$2$ counts on both branches.
\item \emph{Binary lifting.}  The $K$-generic lifting program compiled at $K=13$
  reproduces the output of the audited $k=13$ filter on the $p=239$ rows
  of~\cite{Allikvere26} byte for byte; the $K=14$ build passes a self-test
  of $300$ exact cover cases, the lookups for $\vct{a}$ at three primes,
  and the lift kernel on $4{,}000$ rows from the verified computations.
\item \emph{Level $15$.}  The program agrees with exhaustive enumeration
  for $k=2,\ldots,6$ and reproduces the $k=13$, $p=239$ result.
\item \emph{Repetition.}  $297$ of the $1{,}032$ generation jobs at
  $p=547$ were run twice, on different hosts, with identical node counts
  and rows; $p=233$ was run
  with and without Lemma~\ref{lem:neargcd}, the framework-only run leaving
  $13{,}456$ orbits unresolved, as predicted.
\item \emph{Packages.}  Every package was hash-verified on collection.
\end{enumerate}

\subsection{Limits}
These checks do not constitute a formal proof verification or a fully
independent implementation.  Completeness of generation rests on
Lemmas~\ref{lem:quota} and~\ref{lem:km1all}, together with the verified
precondition where the decomposition variant is used.  Completeness of
the terminal search rests on Lemma~\ref{lem:factor} and the exhaustive
checks described above.  The two accompanying scripts recompute the
analytic constants directly from their definitions.

\section{Computational cost}\label{sec:cost}

\begin{table}[ht]\centering\caption{Cost by stage, in CPU-seconds summed over the jobs of that stage (irredundant branch, reducible branch, precondition, binary lifting, orbit reduction, and level $15$).  Jobs ran in parallel on a sixteen-vCPU machine.  The final column is the total wall time for the prime, in seconds.  For $p=181$ and $p=233$, binary lifting resumed from files written by an interrupted run, so its time was not recorded.}\label{tab:cost}\small
\begin{tabular}{@{}rrrrrrrr@{}}\toprule $p$ & irredundant & reducible & precond. & binary & orbits & level 15 & wall\\ \midrule
89 & 1 & 1 & 0 & 3,514 & 3 & 16,741 & 1,299\\
131 & 32 & 20 & 2 & 148,269 & 54 & 179,778 & 21,504\\
149 & 48 & 12 & 1 & 24,207 & 0 & 77,469 & 7,386\\
181 & 1,206 & 304 & 10 & --- & 252 & 503,678 & ---\\
191 & 512 & 89 & 4 & 63,530 & 0 & 115,568 & 14,303\\
199 & 927 & 119 & 6 & 5,556 & 0 & 9,003 & 1,376\\
233 & 1,862 & 238 & 11 & --- & 0 & 881 & ---\\
239 & 1,301 & 58 & 7 & 15 & 0 & 0.3 & 131\\
241 & 4,566 & 548 & 25 & 683 & 1 & 741 & 1,021\\
307 & 14,558 & 517 & 56 & 10 & 0 & 0.3 & 1,027\\
353 & 24,185 & 746 & 84 & 1 & 0 & 0.3 & 1,668\\
397 & 57,216 & 1,545 & 159 & 1 & 0 & 0.3 & 3,872\\
449 & 70,074 & 1,577 & 173 & 2 & 0 & 0.3 & 4,675\\
499 & 260,613 & 5,611 & 516 & 2 & 0 & 0.3 & 17,237\\
521 & 236,244 & 4,668 & 473 & 2 & 0 & 0.3 & 15,614\\
547 & 444,985 & 8,639 & 774 & 3 & 0 & 0.3 & 29,340\\
557 & 573,547 & 10,928 & 982 & 3 & 0 & 0.3 & 37,681\\
563 & 499,447 & 9,125 & 881 & 3 & 0 & 0.3 & 32,842\\
569 & 351,085 & 6,072 & 737 & 3 & 0 & 0.3 & 23,149\\
\midrule all 71 primes & 5,808,889 & 127,590 & 12,382 & 389,957 & 355 & 1,354,166 & \\
\bottomrule\end{tabular}\end{table}


At the primes with $\tau(p)\geq13$ the irredundant level-$1$ branch
accounts for at least $94\%$ of the CPU time ($p=239$) and for more than
$97\%$ for the verified primes $p>400$; its node count grows like $p^{\CostExponent}$
over the verified set, while binary lifting takes seconds and the
level-$15$ step about a second.  At the primes with $\tau(p)\leq12$,
generation takes minutes, whereas binary lifting and the level-$15$
processing of the $10^4$--$10^7$ surviving unit orbits account for most
of the running time; these tasks run in
parallel.  Each prime was processed on one sixteen-vCPU Hetzner cpx62
machine.  Five machines handled disjoint prime lists, so the wall time was
that of the longest list, about three days.  The total CPU time is
\TotalCpuHours{} hours, or about \TotalCpuYears{} CPU-years.  The machines
supplied no more than $5{,}800$ vCPU-hours of aggregate wall-clock
capacity.  For comparison, the
$LRC(13)$ computation of~\cite{Allikvere26} used $111$ primes with
logarithmic sum $681.5$ and took $27$ days with $31{,}872$ vCPU-hours of
machine capacity.  With the earlier finite-checking threshold, the prime
set would have extended to about $970$; extrapolating the measured growth
gives roughly twenty times the cost of the present computation.

\section{Check of the \texorpdfstring{$LRC(9)$}{LRC(9)} input}\label{app:source}

The chain of lower-dimensional results used in Theorem~\ref{thm:main}
passes through $LRC(9)$.  The program published with
\cite{Trakulthongchai26}, \texttt{lrc\_for\_ten\_runners.cpp} at commit
\texttt{6789649} of~\cite{TrakulthongchaiCode}, forms the intersection of
two lifted families using a key that depends on coordinate order.  Because
the tuples at levels $2p$ and $5p$ are sorted by value rather than by
residue, this step need not form the full intersection.  The author of
\cite{Trakulthongchai26} has been informed and is preparing a corrected
version of the program.

The emptiness of $J(9,p)$ is established independently, without an
intersection step, by the successive-lifting code accompanying~\cite{ST26}
at commit \texttt{755b116} of~\cite{STcode}, in the log
\texttt{results/result\_10}.  That log verifies $J(9,p)=\varnothing$ for
$51$ primes in the range $19\leq p\leq317$.  Their logarithms sum to
$257.87$, which exceeds the finite-checking threshold
$\log B_9=254.30$.  We inspected both programs and checked this
logarithmic sum.  Thus the independent computation accompanying~\cite{ST26}
supplies the $LRC(9)$ input used here.

\begin{thebibliography}{99}
\small
\setlength{\itemsep}{1pt}

\bibitem{Allikvere26}
J. Allikvere,
\emph{Fourteen lonely runners},
\href{https://arxiv.org/abs/2609.02604}{arXiv:2609.02604} (2026);
verification package
\href{https://doi.org/10.5281/zenodo.22066772}{doi:10.5281/zenodo.22066772}.

\bibitem{BCS26}
M. Blanco, F. Criado, and F. Santos,
\emph{Coloopless zonotopes and counterexamples to the Shifted Lonely
Runner Conjecture},
\href{https://arxiv.org/abs/2603.24784}{arXiv:2603.24784} (2026).

\bibitem{GK26}
V. Giri and N. Kravitz,
\emph{The structure of Lonely Runner spectra},
Math.\ Proc.\ Cambridge Philos.\ Soc.\ \textbf{180} (2026), no.~2,
343--361; arXiv:2304.01462.

\bibitem{GW06}
L. Goddyn and E. B. Wong,
\emph{Tight instances of the lonely runner},
Integers \textbf{6} (2006), Paper No.~A38.

\bibitem{KZ1873}
A. Korkine and G. Zolotareff,
\emph{Sur les formes quadratiques},
Math.\ Ann.\ \textbf{6} (1873), 366--389.

\bibitem{LLS90}
J. C. Lagarias, H. W. Lenstra, Jr., and C. P. Schnorr,
\emph{Korkin--Zolotarev bases and successive minima of a lattice and its
reciprocal lattice},
Combinatorica \textbf{10} (1990), no.~4, 333--348.

\bibitem{MSS25}
R. D. Malikiosis, F. Santos, and M. Schymura,
\emph{Linearly exponential checking is enough for the lonely runner
conjecture and some of its variants},
Forum of Mathematics, Sigma \textbf{13} (2025), e164,
\href{https://doi.org/10.1017/fms.2025.10107}{doi:10.1017/fms.2025.10107}.

\bibitem{Rosenfeld25}
M. Rosenfeld,
\emph{The lonely runner conjecture holds for eight runners},
Math.\ Comp.\ (2026), electronically published,
\href{https://doi.org/10.1090/mcom/4243}{doi:10.1090/mcom/4243};
arXiv:2509.14111.

\bibitem{Trakulthongchai26}
T. Trakulthongchai,
\emph{Nine and ten lonely runners},
Electron.\ J.\ Combin.\ \textbf{33} (2026), no.~2, Paper No.~P2.46,
\href{https://doi.org/10.37236/14972}{doi:10.37236/14972}.

\bibitem{TrakulthongchaiCode}
T. Trakulthongchai,
\emph{Nine and ten lonely runners: source code},
GitHub repository,
\href{https://github.com/t-tanupat/nine-and-ten-lonely-runners/commit/6789649}%
{commit \texttt{6789649}} (accessed 9 September 2026).

\bibitem{ST26}
T. Sungkawichai and T. Trakulthongchai,
\emph{Eleven, twelve, and thirteen lonely runners},
\href{https://arxiv.org/abs/2604.23906}{arXiv:2604.23906v2} (2026).

\bibitem{STcode}
T. Sungkawichai,
\emph{LRC holds for 11, 12, 13 runners: source code},
GitHub repository,
\href{https://github.com/vzsky/13-lonely-runners/tree/755b116}%
{commit \texttt{755b116}} (accessed 9 September 2026).

\bibitem{Wills1967}
J. M. Wills,
\emph{Zwei S\"atze \"uber inhomogene diophantische Approximation von
Irrationalzahlen},
Monatshefte f\"ur Mathematik \textbf{71} (1967), no.~3, 263--269,
\href{https://doi.org/10.1007/BF01298332}{doi:10.1007/BF01298332}.

\end{thebibliography}

\end{document}
